\documentclass[11pt,a4paper]{article}
\usepackage[T1]{fontenc}
\usepackage[utf8]{inputenc}
\usepackage{lmodern,amsmath,amssymb,amsthm}
\usepackage[margin=23mm]{geometry}
\usepackage{booktabs,tabularx,array,xcolor,enumitem,microtype}
\usepackage[colorlinks=true,linkcolor=teal!60!black,urlcolor=teal!60!black,citecolor=teal!60!black]{hyperref}
\usepackage{fancyhdr,titlesec,needspace}
\titlespacing*{\section}{0pt}{2.0ex plus .6ex minus .2ex}{.9ex plus .2ex}
\titlespacing*{\subsection}{0pt}{1.5ex plus .4ex minus .2ex}{.7ex plus .2ex}
\pagestyle{fancy}\fancyhf{}
\fancyhead[L]{\small Low-rank Grothendieck rings}
\fancyhead[R]{\small Research announcement}
\fancyfoot[C]{\small\thepage}
\setlength{\headheight}{14pt}
\setlength{\parindent}{0pt}
\setlength{\parskip}{3pt}
\setlist{nosep,leftmargin=*}
\newcommand{\FPdim}{\operatorname{FPdim}}
\newcommand{\Rep}{\operatorname{Rep}}
\newcommand{\Vect}{\operatorname{Vec}}
\newcommand{\Fib}{\operatorname{Fib}}
\newcommand{\NearGroup}{\operatorname{NG}}
\newcommand{\Ad}{\operatorname{Ad}}
\newcommand{\PSU}{\operatorname{PSU}}
\newcommand{\TT}{\mathcal{TT}}
\newcommand{\caseid}[1]{\texttt{\detokenize{#1}}}
\newcommand{\repo}{https://github.com/sebastienpalcoux/grothendieck-ring-classification}
\newcommand{\snap}{https://github.com/sebastienpalcoux/grothendieck-ring-classification/blob/e34396cd95aff1090c10974dd934fe6054ca45eb}
\newcommand{\prooflink}[2]{\href{\snap/#1}{#2}}
\newcommand{\caution}{\textbf{AI-generated research \textperiodcentered\ independent verification pending.}}
\newtheorem{problem}{Problem}
\title{\vspace{-12mm}\textbf{Low-rank Grothendieck rings:}\\[3pt]
an AI-assisted classification project,\\ constructions, and open problems}
\author{S\'ebastien Palcoux\thanks{Contributor and project coordinator. Beijing Institute of Mathematical Sciences and Applications (BIMSA), Beijing, China. AI provenance and verification status are stated in Section~\ref{sec:resources}.}}
\date{17 October 2026}
\hypersetup{pdftitle={Low-rank Grothendieck rings: an AI-assisted classification project, constructions, and open problems},pdfauthor={Sebastien Palcoux},pdfsubject={Research announcement and invitation to independent mathematical review}}
\begin{document}
\maketitle
\vspace{-5mm}
\begin{abstract}
We present a public research project on the Grothendieck rings of complex
fusion categories of rank at most eight. The project reports unrestricted
classification theorems in ranks two, three, and four, and classifications
within specified fusion-multiplicity bounds in ranks five through eight,
with two rank-six categorification questions unresolved. Its resources
include research manuscripts, an exact fusion-ring census, verification
code, and an interactive atlas. We describe seven constructions that are
promising candidates for new tensor categories, subject to a separate
assessment of their novelty up to tensor equivalence. We also state the
two open problems explicitly. This announcement records the project's
claims and their verification status; it invites independent mathematical
review and does not provide a new proof of the complete classification.
\end{abstract}

\begingroup\small
\caution\ The research manuscripts, arguments, code, and data have not yet
been independently human-verified or Lean-formalized. Please treat the
reported results with caution. Exact computations establish only the
mathematical and computational statements within their documented scope.
\endgroup

\section{The classification question and its scope}
Let $\mathcal C$ be a fusion category over $\mathbb C$, with simple-object
classes $b_0=1,b_1,\ldots,b_{r-1}$. Its Grothendieck ring has multiplication
\[
 b_i b_j=\sum_{k=0}^{r-1}N_{ij}^{k}b_k,
 \qquad N_{ij}^{k}\in\mathbb Z_{\geq0}.
\]
Here $r$ is the \emph{rank}, and throughout this announcement
\[
 M(R)=\max_{i,j,k}N_{ij}^{k}
\]
is the \emph{fusion multiplicity}. Thus every bound in Table~\ref{tab:scope}
is a bound on fusion multiplicity, unless explicitly stated otherwise.
A fusion ring is counted up to based-ring isomorphism, preserving the unit
and distinguished basis. To be \emph{realized} means to admit at least one
complex fusion category with that Grothendieck ring. This is a ring-level
classification problem; it does not classify all associators or all tensor
equivalence classes realizing a given ring.

The rank-one case is the ring of $\Vect$. In rank two the familiar family is
$X^2=1+nX$ with $n\geq0$, and only $n=0,1$ admit complex
categorifications~\cite{Ostrik2}. Both have $M=1$, because the unit
coefficients are included in the maximum. The higher-rank account below is
taken from the fixed project snapshot specified in Section~\ref{sec:resources}.

\begin{table}[ht]
\centering\small
\setlength{\tabcolsep}{5pt}
\begin{tabularx}{\textwidth}{c r X r r}
\toprule
Rank & Census through $M$ & Reported general classification & Realized & Open\\
\midrule
2 & all & \textbf{Unrestricted theorem} & 2 & 0\\
3 & $1{,}000$ & \textbf{Unrestricted theorem} & 5 & 0\\
4 & $128$ & \textbf{Unrestricted theorem} & 15 & 0\\
5 & $20$ & Complete through $20$ & 17 & 0\\
6 & $9$ & Complete through $3$; two open through $9$ & 57 & 2\\
7 & $4$ & Complete through $4$ & 29 & 0\\
8 & $2$ & Complete through $2$ & 77 & 0\\
\bottomrule
\end{tabularx}
\caption{Reported status at the pinned snapshot. For ranks five through
eight, realized and open counts refer only to the displayed census range.
The census bounds in ranks three and four describe retained enumeration
data; the reported classification theorems are unrestricted in multiplicity.}
\label{tab:scope}
\end{table}

\section{Reported results and the role of exact computation}
The \prooflink{manuscripts/rank4/papers/companion.pdf}{unrestricted rank-four manuscript}
reports fifteen possible Grothendieck rings, without assuming pivotality or
a multiplicity bound, and states that every categorification admits a
spherical structure. The stated sharp bounds are
\[
 \FPdim(\mathcal C)\leq24+12\sqrt3,\qquad M(K_0(\mathcal C))\leq6.
\]
The dimension bound is attained by the near-group ring $\NearGroup(C_3,6)$.
The proof presented there combines a theoretical reduction to
$\FPdim(\mathcal C)<1000$ and $M\leq31$ with exact finite verification of
$1{,}692$ remaining rings, partitioned into $1{,}677$ exclusions and
$15$ models. The separate census through $M=128$ has $889{,}530$
rings. Related work of Dong, Palcoux, and Plessis treats the pivotal
rank-four case~\cite{DPP}; its hypothesis should be distinguished from
the unrestricted claim of this project.

For ranks five through eight, the complete census totals and recorded
decisions are respectively
\[
\begin{array}{c|r|r|r|r}
 r & \text{census total} & \text{excluded} & \text{realized} & \text{open}\\\hline
5&40{,}999&40{,}982&17&0\\
6&25{,}421&25{,}362&57&2\\
7&4{,}201&4{,}172&29&0\\
8&973&896&77&0
\end{array}
\]
All four partitions concern their stated multiplicity ranges. At rank six,
the only unresolved multiplicities are $4$ and $8$; all $8{,}809$ rings
of exact multiplicity $9$ are recorded as excluded. The $57$ realizations
comprise $56$ unitary models and one spherical model whose ring has no
pseudo-unitary categorification. Thus the target problem is not restricted
to unitary or pseudo-unitary fusion categories.

The approach combines structural constraints on fusion categories with
exact applications to enumerated rings. Its criteria include arithmetic
and cyclotomic restrictions, formal codegrees, induction to the Drinfeld
center, Frobenius--Schur indicators, fusion subcategory and module-category
constraints, and associator or tensor-functor obstructions. Positive
decisions require a categorical model and a comparison of its entire
fusion tensor with the census representative.

These are several distinct verification tasks. A tensor check can establish
associativity of a based ring, a full comparison of structure constants can
identify two based rings, and a finite certificate can establish a specified
obstruction. None of these alone constructs compatible categorical
associators or proves that the census is exhaustive. The manuscripts and
the \prooflink{docs/VERIFY.md}{verification guide} state the relevant
hypotheses, source identifications, and replay domains. The overall project
therefore combines theoretical arguments and exact finite verification.

\Needspace{7\baselineskip}
\section{Seven constructions to examine for novelty}
\textbf{To the best of our current knowledge, these seven constructions
are promising candidates for new tensor categories. Their novelty up to
tensor equivalence remains subject to independent verification.}
They are a selected shortlist of reported constructions from known input
categories, rather than a count of established new tensor equivalence
classes. In particular, identifying a census ring with an already known
category is not itself a discovery of a new category.

In the notation $\caseid{Rr_m_i}$, the first two indices record the rank
and exact maximum fusion multiplicity, and the last is a census identifier.
Write $C_n$ for the cyclic group of order $n$, and $\NearGroup(G,k)$ for a
near-group category with invertible group $G$ and noninvertible simple
$\rho$ satisfying $\rho^2=\sum_{g\in G}g+k\rho$.

\begin{enumerate}[label=\arabic*.,itemsep=5pt]
\item \textbf{\caseid{R6_4_249}: sign-swap equivariantization.}
The reported construction takes the diagonal-fermion quotient of two
identical half-$E_6$ categories and then equivariantizes a factor-swap
action with a sign on the condensed algebra. The known fermionic inputs
are discussed in~\cite{ALW}. See the
\prooflink{manuscripts/rank6/papers/chapters/06.tex}{rank-six construction chapter},
theorem labelled \texttt{thm:sign-swap}.

\item \textbf{\caseid{R6_4_322}: mixed diagonal-fermion quotient.}
Let $\mathcal E$ denote the half-$E_6$ category. The construction uses
the diagonal central boson in $\mathcal E\boxtimes\PSU(2)_6$ and the
corresponding module category. Both factors are known inputs~\cite{ALW}.
See the same chapter,
proposition labelled \texttt{prop:mixed-fermion}.

\item \textbf{\caseid{R6_6_1712}: a relative Fibonacci center.}
The proposed identification is $Z_{\Fib}(\TT_3)$, where the known category
$\TT_3$ is the dual even part of the $2D2$ subfactor~\cite{EM}.
The relative-center framework is standard~\cite{GNN}.
See the same chapter, corollary labelled \texttt{cor:tt3-relative}.

\item \textbf{\caseid{R7_4_35}: an inversion and Morita construction.}
Starting from the $\NearGroup(C_4\times C_4,16)$ input of Yu--Zhang~\cite{YZ},
the project constructs an orbifold, descends a coherent inversion action,
and takes the Morita dual associated to the reflection algebra.
Its four invertibles form $C_2\times C_2$, and all simples are self-dual.
See the \prooflink{evidence/unified_20261016/rank7/R35_INVERSION_CONSTRUCTION.md}{construction note}.

\item \textbf{\caseid{R7_4_2665}: the cyclic near-group analogue.}
The analogous construction starts from Budinski's
$\NearGroup(C_{16},16)$ input~\cite{Budinski}, using the near-group orbifold
and Morita machinery of~\cite{Izumi}. The invertibles form $C_4$;
the two large noninvertibles are mutually dual. Both rank-seven cases have
simple dimensions $(1,1,1,1,2,4+2\sqrt5,4+2\sqrt5)$.
See the \prooflink{evidence/unified_20261016/rank7/R2665_INVERSION_CONSTRUCTION.md}{construction note}.

\item \textbf{\caseid{R8_2_37}: a nonsplit $C_2$-extension of $2D2$.}
Using the constructions in~\cite{Grossman,GIS}, the project transports an
odd automorphism through a crossed product. The odd invertible $a$
satisfies $a^2=g$, where $g$ is the nontrivial invertible in the neutral
$2D2$ category. The simple dimensions are four copies of $1$ and four
copies of $2+\sqrt5$. See the
\prooflink{evidence/unified_20261016/rank8/R37_CONSTRUCTION.md}{construction note}.

\item \textbf{\caseid{R8_2_575}: a self-dual Hopf extension.}
An explicit symmetric Hopf pairing on Kashina's $16$-dimensional
self-dual Hopf algebra~\cite{Kashina} is used to apply the
Davydov--Runkel extension construction~\cite{DR}.
The simple dimensions are $(1,1,1,1,2,2,2,4)$, and the global
Frobenius--Perron dimension is $32$. See the
\prooflink{evidence/upper_residual_20261002b/r8_integral/r575/PROOF.md}{construction note}.
\end{enumerate}

The most useful checks here concern both the existence of the stated
coherence data and comparison with the prior literature. Known input
categories and general extension theorems do not by themselves settle
whether the resulting examples have appeared before.

\section{The two unresolved categorification problems}
The following two based rings remain open in the stated project range.
They have resisted repeated attempts at both categorical construction
and exclusion. The displayed presentations specify concrete targets;
the existence of their fusion tensors does not imply categorifiability.

\subsection*{\caseid{R6_4_155}: a generalized near-group ring}
The commutative based ring has basis $1,h,g,hg,X,Y$, with
\begin{align*}
 g^2=h^2&=1, & gh&=hg, & hX&=X, & gX&=Y=X^*,\\
 XY&=1+h+4X+4Y, & X^2&=g+hg+4X+4Y.
\end{align*}
The remaining products follow from these relations. Its four invertibles
form $C_2\times C_2$, the two noninvertibles are mutually dual, and
\[
 \FPdim(X)=\FPdim(Y)=4+3\sqrt2,
 \qquad \FPdim(R)=72+48\sqrt2.
\]
Its maximum fusion multiplicity is $4$. The ring is already the explicit
candidate $R'(4)$ in Meng--Yu~\cite[Theorem 1.1]{MY}; it is not claimed
here as a newly discovered based ring.

\begin{problem}
Does $\caseid{R6_4_155}$ admit a fusion-category realization over
$\mathbb C$? Construct compatible categorical data or prove that no such
category exists.
\end{problem}

The \prooflink{supplements/rank6/research/frontier/R6_4_155/README.md}{current frontier account}
records a necessary rank-$52$ center profile with a semion factor and
a rank-$26$ centralizer candidate. It also records a rank-$22$ based
module satisfying the stated fermion Azumaya linking identities.
These are necessary-condition calculations: they supply neither a
coherent module category nor associators for a categorification.

\subsection*{\caseid{R6_8_1660}: an Ising-based ring}
The commutative based ring has basis $1,g,X,gX,\sigma,T$, where
$T=\sigma X$, and relations
\[
 g^2=1,\qquad g\sigma=\sigma,\qquad \sigma^2=1+g,
 \qquad X^2=1+\bigl(3(1+g)+4\sigma\bigr)X.
\]
All basis elements are self-dual, and $\{1,g,\sigma\}$ is an Ising
fusion subring. Writing
\[
 d=3+2\sqrt2+2\sqrt3+\sqrt6,
\]
the simple dimensions, in the displayed basis, are
$(1,1,d,d,\sqrt2,\sqrt2d)$ and $\FPdim(R)=4(1+d^2)$.
Its maximum fusion multiplicity is $8$, attained in $T^2$.

\begin{problem}
Does $\caseid{R6_8_1660}$ admit a fusion-category realization over
$\mathbb C$? Determine whether the necessary center and module-category
data can be completed to an actual category.
\end{problem}

The \prooflink{supplements/rank6/research/frontier/R6_8_1660/README.md}{current frontier account}
records relative-center, Morita, and indicator restrictions. For the
positive spherical representative, a bounded necessary-condition domain
contains $640$ labelled restriction profiles;
this is not a list of categories or associators. A particular rank-$37$
candidate fails Verlinde integrality, while other profiles remain open.
Neither this rejection nor the survival of other profiles decides the
general categorification question.

\section{Resources, provenance, and independent review}\label{sec:resources}
The project has three complementary public entry points:
\begin{itemize}[itemsep=3pt]
\item \textbf{Classification manuscripts, models, and decision records:}\\
\url{https://github.com/sebastienpalcoux/grothendieck-ring-classification}.
\item \textbf{Fusion-ring census and enumeration sources:}\\
\url{https://github.com/sebastienpalcoux/fusion-ring-census}.
\item \textbf{Interactive atlas:}\\
\url{https://sebastienpalcoux.github.io/grothendieck-ring-classification/}.
\end{itemize}
The \prooflink{docs/New_category_candidates_and_open_cases.pdf}{one-page construction shortlist}
provides a quick entry into Section~3. The
\prooflink{data/remaining_cases.py}{two exact open fusion tensors}
are available as Python lists with convention $N[i][j][k]=N_{ij}^{k}$,
zero-based unit index $0$, explicit duality maps, and tensor checksums.
The symbolic basis orders in Section~4 need not coincide with the
canonical file order; the first case is obtained from canonical indices
$(0,3,4,5,1,2)$, and the second already has the displayed order.

\textbf{Fixed research snapshot.}
This announcement describes classification commit
\href{https://github.com/sebastienpalcoux/grothendieck-ring-classification/tree/e34396cd95aff1090c10974dd934fe6054ca45eb}{\texttt{e34396cd95af}}
and census-repository commit
\href{https://github.com/sebastienpalcoux/fusion-ring-census/tree/d4c36cc2ee5eb5f0e5291de18d2bb433ffb7ce8b}{\texttt{d4c36cc2ee5e}}.
Rank-specific source manifests retain the exact census inputs used for
individual certificates; the displayed census-repository snapshot does not
replace those source bindings. Live pages may evolve, whereas these
commit links identify the resources summarized here.

The research collection at this snapshot comprises eight classification
manuscripts and six census manuscripts: $14$ research manuscripts and
$485$ pages in total, including references and appendices. Verification
guides and this announcement are separate documents. The
\prooflink{manuscripts/README.md}{classification reading guide} identifies
the unrestricted rank-four paper as the primary rank-four account.

\textbf{AI provenance.} The project was developed through AI-assisted
research under the contributor's direction; its research text, proposed
arguments, code, and data are primarily AI-generated. This announcement
was also prepared with ChatGPT. Automated rechecking and AI-assisted
revision have been used, but do not constitute independent human
verification. No Lean formalization of the full project is claimed.
Historical novelty of the shortlisted categorical constructions is a
separate question from their mathematical validity. \caution

Readers can select a focused task: verify a stated categorical construction
and its coherence data; compare a model with known tensor equivalence
classes; inspect a classification reduction and its finite certificate
scope; or construct or exclude one of the two open rings. For code and
data, the \prooflink{docs/VERIFY.md}{replay guide} and
\prooflink{docs/PROVENANCE.md}{source-provenance record} identify the
relevant inputs and exact checks. We welcome comparisons with previously
known tensor categories, independent scrutiny of the proposed
constructions and classification arguments, and progress on the two
remaining categorification problems.

\begin{thebibliography}{99}\small
\setlength{\itemsep}{1pt}\setlength{\parskip}{0pt}
\bibitem{Ostrik2} V. Ostrik, \emph{Fusion categories of rank 2},
Math. Res. Lett. \textbf{10} (2003), 177--183.
\href{https://arxiv.org/abs/math/0203255}{arXiv:math/0203255}.
\bibitem{DPP} J. Dong, S. Palcoux, and A. Plessis,
\emph{Classifying Grothendieck rings of pivotal fusion categories of rank four},
\href{https://arxiv.org/abs/2609.33272}{arXiv:2609.33272} (2026).
\bibitem{ALW} D. Aasen, E. Lake, and K. Walker,
\emph{Fermion condensation and super pivotal categories},
\href{https://arxiv.org/abs/1709.01941}{arXiv:1709.01941} (2017; revised 2018).
\bibitem{EM} C. Edie-Michell,
\emph{A complete classification of unitary fusion categories tensor generated
by an object of dimension $(1+\sqrt5)/2$},
\href{https://arxiv.org/abs/1904.08909}{arXiv:1904.08909} (2019; revised 2020).
\bibitem{GNN} S. Gelaki, D. Naidu, and D. Nikshych,
\emph{Centers of graded fusion categories},
\href{https://arxiv.org/abs/0905.3117}{arXiv:0905.3117} (2009).
\bibitem{YZ} Z. Yu and Q. Zhang,
\emph{Realizing modular data from centers of near-group categories},
\href{https://arxiv.org/abs/2412.20763}{arXiv:2412.20763} (2024; revised 2025).
\bibitem{Budinski} P. G. Budinski,
\emph{Exotic Fusion Categories and Their Modular Data},
master's thesis, University of Alberta, 2021.
\href{https://doi.org/10.7939/r3-7p7k-fg18}{doi:10.7939/r3-7p7k-fg18}.
\bibitem{Izumi} M. Izumi, \emph{A Cuntz algebra approach to the
classification of near-group categories},
\href{https://arxiv.org/abs/1512.04288}{arXiv:1512.04288} (2015).
\bibitem{Grossman} P. Grossman,
\emph{Fusion categories associated to subfactors with index $3+\sqrt5$},
\href{https://arxiv.org/abs/1612.07185}{arXiv:1612.07185} (2016).
\bibitem{GIS} P. Grossman, M. Izumi, and N. Snyder,
\emph{Graded extensions of generalized Haagerup categories},
\href{https://arxiv.org/abs/2201.11901}{arXiv:2201.11901} (2022).
\bibitem{Kashina} Y. Kashina,
\emph{Classification of semisimple Hopf algebras of dimension 16},
\href{https://arxiv.org/abs/math/0004114}{arXiv:math/0004114} (2000).
\bibitem{DR} A. Davydov and I. Runkel,
\emph{$\mathbb Z/2\mathbb Z$-extensions of Hopf algebra module categories by their base categories},
\href{https://arxiv.org/abs/1207.3611}{arXiv:1207.3611} (2012; revised 2013).
\bibitem{MY} Y. Meng and Z. Yu,
\emph{Classification of some $\mathbb Z/2\mathbb Z\times\mathbb Z/2\mathbb Z$-quadratic
fusion categories of rank 6},
\href{https://arxiv.org/abs/2608.06064}{arXiv:2608.06064} (2026).
\end{thebibliography}
\end{document}
